\section{Related Work}\label{related-work}

\subsection{Landauer, Non-Equilibrium Thermodynamics, and Selection}

Landauer's principle gives the standard calibration from logically irreversible discrimination to minimum heat production and energy cost \cite{landauer1961irreversibility,bennett1982thermodynamics}. Stochastic thermodynamics extends that floor to trajectory-level entropy production, work identities, and fluctuation relations \cite{seifert2012stochastic,vandenbroeck2015ensemble,wolpert2019stochastic,jarzynski1997nonequilibrium,crooks1999entropy}. Finite-time erasure and mismatch corrections sharpen the same theme for controlled nonequilibrium protocols \cite{diana2013finite,proesmans2020finite,manzano2024absolute}. The theorem chain above isolates a different object: a finite exact-resolution lower bound indexed by the number of independent coordinates that must be resolved to preserve the optimizer.

Relative to the Seifert and Van den Broeck--Esposito framework, the present model does not attempt a full trajectory description of a driven Markov process, housekeeping heat, or protocol-dependent dissipation. It gives instead a mechanized non-asymptotic lower bound in terms of structural rank and decision-quotient entropy. Conversely, stochastic-thermodynamic frameworks resolve time-dependent nonequilibrium refinements that are outside the current finite exact-resolution model.

Non-equilibrium selection and replication arguments lie in the same neighborhood \cite{england2013statistical}. The entropy premium used above is deliberately discrete: multiplicity is reduced to counting over a finite quotient family and the elementary inequality $k \le 2^{k-1}$, rather than to a separate stochastic-process ansatz.

\subsection{Zero-Error, Functional, and Quotient Information}

The information object is closer to zero-error and confusability-based information theory than to average-case source coding \cite{shannon1956zero,korner1973graphs,lovasz1979shannon,csiszar2011information}. The central quantity is not full state entropy but the entropy of the decision quotient: how many distinct optimal-action classes survive after irrelevant coordinates are erased.

Function-relative information in physics and origins-of-life work also conditions information on successful function or selection \cite{szostak2003functional,wong2023roles}. The object studied here is narrower and exact: coordinate erasure is admissible precisely when optimal-action correspondence is preserved. The rank-$1$ regime is therefore the one-coordinate exact-decision regime, the tractable sufficiency regime, and the minimum calibrated-cost regime.

\subsection{Categorical Quotients and Exact Abstraction}

Quotienting states by equality of $\Opt$ is the standard coimage construction for the decision quotient of the optimizer map $\Opt : S \to \mathcal{P}(A)$ in \textbf{Set}, canonically equivalent to its image \cite{maclane1998categories}. The novelty is not the existence of that quotient, but the theorem chain tying it to coordinate sufficiency, structural rank, decision entropy, and thermodynamic cost in one proof object.

\subsection{Formal Verification and Machine-Checked Theory}

The proofs are machine-checked in Lean~4 \cite{moura2021lean4} against the Mathlib library \cite{mathlib2020}. Related mechanized precedents include verified computability and semantics developments in Coq and Isabelle \cite{forster2019verified,nipkow2002isabelle,nipkow2014concrete} and certificate-carrying proof artifacts \cite{necula1997proof}. Mechanization matters because the main identifications are exact rather than rhetorical: degree of freedom, structural rank, quotient entropy, tractable sufficiency, and calibrated thermodynamic cost are kept distinct and then linked by explicit theorems.
